\section{Evaluation Method}
\label{sec:method}

Three programs ran back to back on the same product (\cref{tab:programs}). V1 tested team designs
and efficiency on a suite we built; V8 tested levers on Terminal-Bench~4.0 through Harbor; a final
run put the chosen configuration on Terminal-Bench~4.1 the way Artificial Analysis runs agents.

\begin{table}[t]
\centering
\small
\caption{The three programs. Trials in V8 count every graded attempt, including screens and
re-runs.}
\label{tab:programs}
\begin{tabularx}{\textwidth}{@{}l l L >{\raggedright\arraybackslash}p{30mm} >{\raggedright\arraybackslash}p{33mm}@{}}
\toprule
Program & Dates (2026) & Tasks & Models & Scale \\
\midrule
V1 & 28 Sep -- 26 Oct & Own suite: authored, repository, public (Terminal-Bench, DeepSWE) and
  question-answering tasks with hidden graders & GPT-6 Astra, then GPT-6.1 Sol at medium effort &
  Over 120 batches on two local evaluation instances \\
V8 & 27 Oct -- 19 Nov & Terminal-Bench~4.0: the 63 tasks that run without a GPU, split 33 readable
  and 30 unread & GPT-6.1 Sol, GPT-6.1 Terra (screens), GPT-6.2 Sol & 315 Harbor jobs, 6,223
  graded trials on one Linux PC \\
Final & 20 Nov & Terminal-Bench~4.1 at AA's dataset digest, 63 tasks & GPT-6.2 Sol at max effort &
  189 attempts, 3 per task \\
\bottomrule
\end{tabularx}
\end{table}

\subsection{Task sets}

\para{The V1 suite} Each task ships a hidden grader and a reference solution. Before a task entered
the suite, a fresh model solved it blind from an empty workspace, and its grade checked the
grader; this found two real grader defects in one early batch. Saturated or unfair tasks were
retired (52 on 9 October, leaving 32 active). A design that won on the tasks it was tuned on had to
hold on tasks it was not tuned on before it shipped.

\para{The V8 split} On 30 October the 63 runnable Terminal-Bench~4.0 tasks were split into 25
development and 38 held-out tasks. After the first held-out read (4 November), the 8 held-out tasks
on which a board agent beat us were opened for study. That left 33 \emph{readable} tasks, whose
trajectories, failing tests and task files we read to design levers, and 30 \emph{unread} tasks for
confirmation. The unread tasks were seen only as summed rewards, costs, minutes and model calls; no
name, trajectory or test of theirs was opened, then or for this paper.

\subsection{Infrastructure}

V8 ran Harbor 0.24.0 (0.26.0 for the final run) \citep{harbor2026} on a desktop PC with 32 hardware
threads and 96\,GB of memory, about 30 trials at a time after a swap storm at higher load lost 44
runs. The team's own time limit was 465 minutes against the tasks' 480, because the server's
3-hour default had cut an 8-hour task short. Before 5 November, runs used a Docker virtual machine
on a Mac at two trials at a time; its held-out baseline was lower than the PC's, one reason every
comparison is made within one batch on one host.

\subsection{Metrics}

A trial's reward $r_{t,d}$ for task $t$ and draw $d$ is 0 or 1. Our score is the sum over tasks of
each task's mean reward,
\begin{equation}
  S = \sum_{t \in T} \frac{1}{n_t} \sum_{d=1}^{n_t} r_{t,d},
  \label{eq:score}
\end{equation}
so $S/|T|$ is the pass rate. Missing or reward-less attempts count as 0. Cost is the run's tokens
priced at OpenAI's API list price from the server's price table (checked 18 October 2026),
although most runs used a ChatGPT (Codex) sign-in. Time is the agent's own time per attempt; V1's
efficiency work used active time up to the last model call.

\subsection{Statistical analysis}

On the Terminal-Bench sets, an arm and its base each ran two or more draws per task. We tested the
difference in $S$ by re-dealing draws between arm and base within each task (20,000
permutations), which gives the null standard deviation and a $p$-value. With two draws a
side on the unread 30, the null standard deviation is about 1.4 (1.33 to 1.53 across reads), so one
read resolves differences of about 2.7 tasks; with four draws a side it is about 0.94. Where a
standard error is given, it is over tasks. V1 compared paired runs with exact Wilcoxon signed-rank
tests on log ratios of cost or time \citep{wilcoxon1945}.

\subsection{Experimental hygiene}

\begin{itemize}[leftmargin=1.4em]
  \item \emph{Same-batch bases.} Every arm ran beside its base in the same hours, because one base
    moved by +1.67 between batches.
  \item \emph{Quiet re-grades.} Every zero was re-graded once on an idle machine, because verifier
    timeouts under load produced false zeros. On an idle machine they were rare: in two large
    re-grades 43 of 43 and 54 of 54 zeros stayed zero, and in a third 1 of 74 came back.
  \item \emph{Registration.} The rule, the branches and our own forecast were written into a lab
    notebook before each batch \citep{nosek2018prereg}; the V8 notebook reached 9,276 lines in 387
    sections. An independent reviewer agent was asked to find flaws in designs and reads.
\end{itemize}
